\documentclass[3p,times]{elsarticle}

% ===============================================================
% PACKAGES
% ===============================================================

\usepackage[utf8]{inputenc}

\usepackage{multirow}
\usepackage{amsmath,amssymb,amsfonts}
\usepackage{amsthm}
\usepackage{mathrsfs}
\usepackage{xcolor}
\usepackage{textcomp}
\usepackage{manyfoot}
\usepackage{booktabs}

\usepackage{algorithm}
\usepackage{algorithmicx}
\usepackage{algpseudocode}

\usepackage{listings}
\usepackage{url}
\usepackage{hyperref}

% ===============================================================
% HYPERREF SETUP
% ===============================================================

\hypersetup{
    colorlinks=true,
    linkcolor=blue,
    citecolor=blue,
    urlcolor=blue
}

% ===============================================================
% THEOREM ENVIRONMENTS
% ===============================================================

\theoremstyle{plain}
\newtheorem{theorem}{Theorem}

\newtheorem{proposition}[theorem]{Proposition}

\theoremstyle{definition}
\newtheorem{example}{Example}
\newtheorem{remark}{Remark}
\newtheorem{definition}{Definition}

% ===============================================================
% DOCUMENT SETTINGS
% ===============================================================

\raggedbottom

% ===============================================================
% DOCUMENT
% ===============================================================

\begin{document}

\begin{frontmatter}

% ===============================================================
% TITLE AND AUTHORS
% ===============================================================

\title{FairLight: Fairness-Aware City-Wide Emergency Vehicle
Prioritization via Hierarchical Multi-Agent Reinforcement Learning}

\author{Dr G Poonthalir\corref{cor1}}
\author{A Sree Ram}
\author{Aswajith S}

\cortext[cor1]{Corresponding author}

\affiliation{organization={},
            addressline={},
            city={},
            postcode={},
            state={},
            country={}}

% ===============================================================
% ABSTRACT
% ===============================================================

\begin{abstract}

Urban traffic signal control is a complex sequential decision-making
problem in which decisions made at one intersection can influence
traffic conditions at neighboring and downstream intersections. The
problem becomes increasingly challenging as the size of the road
network grows because a controller must simultaneously achieve
network-level coordination, local responsiveness, and computational
scalability. Reinforcement learning has emerged as a promising approach
for adaptive traffic signal control, while multi-agent reinforcement
learning (MARL) allows control to be distributed among individual
intersections. However, fully decentralized approaches may lack
sufficient global context, whereas centralized approaches become
difficult to scale.

Hierarchical MARL frameworks such as HiLight reduce this problem to a
Meta-Policy/Sub-Policy hierarchy coordinated over a fixed, manually
preprocessed region partition and a constant per-intersection
neighborhood, and do not model emergency-vehicle (EMV) priority or the
fairness cost that priority imposes on ordinary traffic. This paper
presents FairLight, a Hierarchical Region-Aware Multi-Agent
Reinforcement Learning framework that (i) replaces fixed regional
preprocessing with automatic, graph-based partitioning --- Louvain and
Leiden community detection over a traffic-weighted intersection graph,
spatial DBSCAN, and four hybrid combinations of the two --- evaluated
as an experimental variable rather than assumed fixed; (ii) augments
each intersection's local observation with a graph-attention
neighborhood representation capped at four neighbors (truncated or
zero-padded as needed, not an arbitrary-degree mechanism); and (iii)
adds an explicit fairness penalty to the Sub-Policy reward, combining
per-intersection waiting-time variance, a starvation threshold, and
Jain's Fairness Index, so that aggregate efficiency gains cannot be
reported without also reporting their distributional cost to ordinary
vehicles.

The framework is implemented and trained with Proximal Policy
Optimization and Generalized Advantage Estimation on top of a
Transformer/LSTM Meta-Policy and a graph-attention Sub-Policy, and is
configured for five SUMO traffic networks (a 472-intersection Manhattan
extract, the Cologne8, Ingolstadt21, Arterial4x4, and Grid4x4 benchmark
networks). City-wide coordination of multiple simultaneously active
EMVs with conflict arbitration, and incident/accident-robust training,
are presented as proposed extensions to this architecture (Sections
3.2 and 3.3): no emergency vehicles are present in any of the five
scenarios' current demand data, and no corresponding training or
evaluation code has been implemented; a formal envy-free objective is
likewise not implemented, and the fairness mechanism above should not
be read as an envy-free guarantee. What has been measured is the
automatic regionalization step itself: applied to the five scenarios,
Louvain and Leiden produce between 1 and 70 regions depending on
network topology, with modularity ranging from near zero on small,
sparsely connected networks to above 0.6 on Manhattan, while DBSCAN's
purely spatial clustering yields coarser, near-uniform partitions
across all five. Full closed-loop evaluation of the trained controller
against Fixed-Time and Max-Pressure baselines has not yet been
completed at the time of writing; the corresponding results are
reported as pending rather than estimated. This manuscript is
accordingly a methodological and partially evaluated contribution: it
documents what is implemented and measured, and separates these
findings explicitly from what remains proposed future work.

\end{abstract}

\begin{keyword}
Traffic Signal Control \sep
Multi-Agent Reinforcement Learning \sep
Emergency Vehicle Routing \sep
Community Detection \sep
Louvain Algorithm \sep
Fairness-Aware Optimization \sep
Hierarchical Reinforcement Learning
\end{keyword}

\end{frontmatter}

% ===============================================================
% INTRODUCTION
% ===============================================================

\section{Introduction}\label{sec1}

Urban traffic signal control is a complex sequential decision-making
problem where decisions at one intersection propagate to neighboring
and downstream intersections. Traffic demand varies dynamically with
time, incidents, weather, and travel patterns, making static or
rule-based controllers insufficient for real-world conditions.
Adaptive systems such as SCOOT introduced traffic-responsive
coordination \cite{scoot}, while reinforcement learning approaches
have since enabled data-driven signal control
\cite{intellight}. Multi-agent reinforcement learning (MARL) further
distributes control across intersections, improving scalability,
though fully decentralized approaches risk optimizing locally at the
expense of network-wide coordination \cite{thousandlights,ma2c}.

Beyond aggregate congestion, traffic signal decisions simultaneously
affect travel time, delay, queue formation, throughput, fuel
consumption, and emissions. Critically, evaluating controllers solely
through aggregate metrics such as average waiting time conceals how
service is distributed among individual vehicles. Consider two
controllers with vehicle waiting times
$W_A = [10, 10, 10, 10, 10, 100]$ and
$W_B = [25, 25, 25, 25, 25, 25]$: Controller~A achieves a lower mean
delay yet inflicts severe starvation on one vehicle. This motivates
fairness as a first-class control objective rather than a post-hoc
evaluation criterion.

A further dimension of real-world uncertainty that existing MARL
frameworks largely overlook is the presence of emergency vehicles
(EMVs). Most traffic signal control approaches, including hierarchical
methods such as HiLight~\cite{hilight} and cooperative approaches such
as CoLight~\cite{colight}, assume homogeneous vehicle populations and
optimize purely for aggregate efficiency. In practice, ambulances,
fire engines, and police vehicles impose stochastic, high-priority
demands on the network that cannot be resolved by ordinary signal
timing. Delayed EMV passage has measurable consequences: studies
report approximately 700 annual fatalities attributable to late
ambulance responses, and an estimated annual mean of 4{,}500
ambulance-involved crashes in the United States
\cite{nellore2016}. EMVLight demonstrated that jointly optimizing EMV
routing and adaptive signal control within a multi-agent actor-critic
framework substantially reduces emergency response times
\cite{emvlight}, but does not address hierarchical city-wide
coordination or fairness. These considerations establish EMV
responsiveness as a necessary component of any realistic traffic
management framework, not an optional extension. At the same time,
providing unrestricted priority to an EMV may excessively delay
ordinary traffic, potentially transferring congestion to other parts
of the network rather than resolving it. EMV prioritization should
therefore not be treated as an isolated intersection-level problem: it
requires coordination across the surrounding traffic network while
maintaining acceptable service for non-emergency vehicles.

Hierarchical reinforcement learning addresses the
scalability--coordination tradeoff by separating long-horizon
strategic decisions from short-horizon local control
\cite{feudalhrl,hirol}. HiLight is particularly relevant here,
employing a Meta-Policy with Transformer-based spatial and LSTM-based
temporal modelling to coordinate decentralized Sub-Policy agents
across large networks~\cite{hilight}. As reported by its authors, its
region partitioning is a preprocessing step rather than a
control-relevant variable, and it does not incorporate EMV priority
or vehicle-level fairness; we treat this characterization of HiLight as
the motivating baseline throughout this paper rather than as a
re-verified, independently reproduced claim.

This paper presents FairLight, a Hierarchical Multi-Agent Actor--Critic
framework that investigates three extensions to this class of
baseline: (1) replacing a fixed or manually specified region partition
with automatic, graph-based regionalization, evaluated as an
experimental variable rather than a fixed preprocessing choice; (2) an
explicit, measured fairness penalty on ordinary-vehicle waiting-time
disparity, in place of a purely aggregate efficiency objective; and (3)
a proposed design, not yet implemented, for coordinating multiple
simultaneously active emergency vehicles with conflicting routes,
motivated by the observation that single-EMV preemption strategies
degrade ordinary-traffic efficiency and do not extend trivially to
concurrent EMVs. The paper's central research questions are: (RQ1)
does the choice of automatic regionalization method materially change
the region structure handed to the Meta-Policy, and by how much, across
networks of different size and topology; (RQ2) can an explicit
fairness penalty be added to the Sub-Policy reward without requiring a
separate architecture; and (RQ3) what would be required, beyond the
existing architecture, to support city-wide multi-EMV coordination.
RQ1 is answered empirically in Section~\ref{sec6} using five SUMO
networks; RQ2 is answered at the level of reward design and
implementation (Section~\ref{sec4}), with closed-loop evaluation
pending; RQ3 is addressed as a design proposal (Section~\ref{subsec3-2})
rather than a validated result, since no part of the multi-EMV
mechanism is currently implemented and none of the five evaluation
scenarios currently include an emergency vehicle in their demand data.

The contributions of this paper, each labeled by its evidentiary
status, are:
\begin{enumerate}
\item \textbf{Implemented and measured.} Automatic regionalization of
the signalized road network via Louvain and Leiden community detection
over a traffic-weighted intersection graph, spatial DBSCAN, and four
hybrid compositions of the two, with region counts, cluster sizes, and
modularity reported for five networks (Section~\ref{sec6}).
\item \textbf{Implemented, evaluation pending.} A hierarchical
Meta-Policy/Sub-Policy architecture trained end-to-end with Proximal
Policy Optimization, in which the Sub-Policy reward includes an
explicit fairness penalty (waiting-time variance, a starvation
threshold, and Jain's Fairness Index) in addition to local traffic
efficiency (Sections~\ref{sec3}--\ref{sec4}).
\item \textbf{Implemented, scope-limited.} A graph-attention
neighborhood mechanism that incorporates information from up to four
neighboring intersections per decision step, with truncation or
zero-padding outside that range (Section~\ref{subsec3-7}).
\item \textbf{Proposed, not implemented.} A city-wide multi-EMV
routing and signal-priority arbitration mechanism for simultaneously
active emergency vehicles (Section~\ref{subsec3-2}), and an
incident/accident-robustness training mechanism (Section~\ref{subsec3-3}).
\end{enumerate}
To the best of our knowledge, no existing framework combines automatic,
experimentally-compared regional partitioning with an explicit
ordinary-traffic fairness penalty inside a hierarchical MARL
controller; the multi-EMV and formal fairness-guarantee extensions
remain future work rather than demonstrated contributions.


% ===============================================================
% RELATED WORK
% ===============================================================

\section{Related Work}\label{sec2}

Conventional traffic signal control has primarily relied on analytical
optimization and rule-based methods for determining optimal signal
timings. Webster's signal timing formulation introduced one of the
earliest mathematical models for fixed-cycle traffic signal
optimization, while later adaptive traffic control systems such as
SCOOT and SCATS dynamically adjusted signal timings using detector-based
traffic measurements and predefined control strategies
\cite{webster1958,scoot,scats,opac}. Max-pressure control further
improved network throughput by selecting signal phases according to
queue pressure differences between incoming and outgoing links
\cite{maxpressure}. Although these approaches are computationally
efficient and have been successfully deployed in practice, they rely
heavily on predefined traffic models and handcrafted optimization
rules. Consequently, they struggle to adapt to highly dynamic urban
environments involving unpredictable congestion, heterogeneous
intersections, accidents, and emergency vehicle requests.

To improve adaptability, several optimization and machine learning
techniques have been introduced for traffic signal control.
Evolutionary optimization methods such as Genetic Algorithms optimize
signal timing plans through iterative search, while Fuzzy Logic
Controllers employ expert-defined rules to respond to uncertain
traffic conditions \cite{genetic,fuzzy}. Later studies adopted machine
learning techniques including Support Vector Machines, Artificial
Neural Networks, and Long Short-Term Memory (LSTM) networks to predict
traffic flow and congestion using historical observations
\cite{anntraffic,svmtraffic,lstmtraffic}. Although these methods
improve prediction accuracy and optimization quality, they primarily
estimate future traffic states rather than learning sequential traffic
control policies. Furthermore, they require extensive historical
training data and cannot continuously adapt to rapidly changing
traffic conditions.

Recent advances in deep reinforcement learning have significantly
improved adaptive traffic signal control by enabling traffic
controllers to learn directly through interaction with the traffic
environment. IntelliLight demonstrated the effectiveness of Deep
Q-Networks for adaptive traffic signal control at isolated
intersections \cite{intellight}. PressLight integrated max-pressure
traffic theory into deep reinforcement learning to maximize traffic
throughput while reducing handcrafted reward engineering
\cite{presslight}. FRAP modeled traffic control as a phase competition
problem and exploited traffic phase symmetry to improve learning
efficiency \cite{frap}. AttendLight later introduced an
attention-based architecture capable of generalizing across
heterogeneous intersections with varying lane configurations
\cite{attendlight}. MetaLight focused on transferring learned
policies across unseen intersections \cite{metalight}, while
Advanced-XLight demonstrated that effective traffic state
representation can significantly influence performance in
reinforcement-learning-based traffic signal control
\cite{advancedxlight}. Despite these advances, existing deep
reinforcement learning methods primarily optimize local traffic
efficiency and do not jointly consider emergency vehicle
prioritization, fairness, or large-scale hierarchical coordination.

To improve cooperation among multiple intersections, graph neural
networks and multi-agent reinforcement learning have received
increasing attention. CoLight introduced graph attention networks to
enable communication among neighboring intersections for cooperative
traffic signal control \cite{colight}. MPLight further improved
scalability by introducing parameter sharing across reinforcement
learning agents, allowing efficient deployment over city-scale
transportation networks \cite{mplight}. HiLight extended this concept
by proposing a hierarchical reinforcement learning framework
consisting of Meta-Policy and Sub-Policy agents for coordinating
thousands of intersections efficiently \cite{hilight}. CityFlow
provides a large-scale multi-agent reinforcement learning environment
designed specifically for city traffic scenarios and supports
evaluation of scalable traffic-control algorithms \cite{cityflow}.
Although these methods substantially improve scalability and
inter-agent cooperation, their primary objective remains maximizing
traffic efficiency. They neither address emergency vehicle routing
nor incorporate fairness-aware resource allocation across different
regions of a city.

HiLight's hierarchical Meta-Policy/Sub-Policy formulation
\cite{hilight}, like most of the methods discussed above, operates
under a closed-world assumption: it is trained and evaluated on
recurrent congestion patterns and does not model non-recurrent
disruptions such as emergency vehicle preemption or traffic
incidents. This omission is not merely conceptual. Poorly coordinated
EMV preemption has been shown to increase average delay for
cross-traffic by 18--32 seconds per vehicle per preemption event, with
these localized delays cascading into network-wide congestion
\cite{dtpreemption}. Even well-established preemption strategies carry
a measurable efficiency cost: EMVLight's own evaluation shows that
traditional preemption techniques increase average non-EMV travel
time by approximately 10\% relative to unpreempted baselines
\cite{emvlight}. These results indicate that emergency vehicle
disruption is a first-order cost of operating a real traffic network,
not a negligible edge case that a closed-world controller can safely
ignore. Motivated by this, our framework treats EMVs and accident
events as first-class constraints on the control problem rather than
external events layered on top of it, at the cost of an increase in
average travel time and delay relative to a disruption-free
benchmark---a cost we treat as the price of a tighter approximation to
real-world operating conditions.

Emergency vehicle routing and traffic signal pre-emption have also
received considerable research attention. Conventional emergency
routing methods generally employ shortest-path algorithms such as
Dijkstra's algorithm, A*, and various time-dependent shortest path
formulations to determine optimal routes under static or partially
dynamic traffic conditions
\cite{dijkstraemv,astar,tdsp,dynamicrouting,miprouting}. More recent
studies have introduced reinforcement learning-based routing
strategies capable of adapting to changing traffic environments
\cite{rlrouting}. EMVLight jointly optimizes emergency vehicle routing
and adaptive traffic signal control using multi-agent actor--critic
reinforcement learning, significantly reducing emergency response
times while maintaining acceptable traffic efficiency for ordinary
vehicles \cite{emvlight}. Subsequently, hierarchical multi-agent
actor--critic methods have been proposed to further improve emergency
routing through coordinated traffic signal control
\cite{hierarchicalemv}. However, existing emergency traffic management
approaches generally assume decentralized coordination and do not
explicitly consider fairness, topology-aware regional decomposition,
or simultaneous optimization for multiple emergency vehicles
operating within large urban transportation networks. Visual
sensing-based approaches have further demonstrated that incorporating
real-time distance measurement and vehicle counting into traffic
management can substantially improve EMV scheduling and intersection
clearance \cite{nellore2016}.

Fairness has recently emerged as an important objective in intelligent
transportation systems. FairMARL introduced a fairness-aware
multi-agent reinforcement learning framework based on generalized
social welfare functions to reduce disparities among intersections
while maintaining competitive traffic efficiency \cite{fairmarl}.
Although fairness-aware optimization improves equitable resource
allocation, existing approaches primarily focus on balancing
intersection delays and do not integrate emergency vehicle routing,
hierarchical coordination, or topology-aware traffic management into
a unified framework.

Overall, significant progress has been achieved in analytical
optimization, deep reinforcement learning, graph-based coordination,
emergency vehicle routing, and fairness-aware traffic signal control.
However, these objectives have largely been investigated independently.
Classical optimization methods lack adaptability, deep reinforcement
learning approaches prioritize traffic efficiency over fairness,
graph-based coordination methods do not explicitly support emergency
response, emergency vehicle routing methods overlook hierarchical
city-wide coordination, and fairness-aware methods ignore emergency
traffic management.

To the best of our knowledge, no existing framework combines
hierarchical multi-agent reinforcement learning with
automatically-compared, graph-based regional partitioning and an
explicit ordinary-traffic fairness penalty within a single city-wide
architecture. FairLight, presented in Section~\ref{sec3}, implements
this combination; city-wide multi-emergency-vehicle routing and a
formal envy-free fairness guarantee are identified as open extensions
rather than claimed as implemented contributions (Section~\ref{sec8}).

Table~\ref{tab:related-comparison} summarizes how FairLight's verified
and proposed components relate to the closest prior systems discussed
above. Entries are marked ``Yes'' only where the cited paper or the
FairLight implementation audited in this work explicitly supports the
capability, ``No'' where the source explicitly does not, and ``Not
established'' where the manuscript cannot verify the capability one way
or the other from the cited source alone.

\begin{table}[t]
\centering
\caption{Comparison of FairLight against the closest related systems.
``Not established'' indicates the capability could not be verified from
the cited source and should not be read as a confirmed absence.}
\label{tab:related-comparison}
\resizebox{\textwidth}{!}{%
\footnotesize
\begin{tabular}{lcccccc}
\toprule
Method & Hierarchical & Automatic & Explicit & Simultaneous & Ordinary-traffic & Formal \\
 & control & regionalization & EMV routing & multi-EMV & fairness & envy-free \\
\midrule
CoLight~\cite{colight} & No & No & No & No & No & No \\
MPLight~\cite{mplight} & No & No & No & No & No & No \\
HiLight~\cite{hilight} & Yes & No (fixed) & No & No & No & No \\
EMVLight~\cite{emvlight} & No & No & Yes & Not established & No & No \\
FairMARL~\cite{fairmarl} & Not established & Not established & No & No & Yes & Not established \\
FairLight (this work) & Yes & Yes (compared) & Proposed & Proposed & Yes & Not implemented \\
\bottomrule
\end{tabular}%
}
\end{table}


% ===============================================================
% PROPOSED FRAMEWORK
% ===============================================================

\section{Proposed Framework}\label{sec3}

The proposed framework, FairLight, is organized around three
components, two of which are implemented and evaluated in this paper
and one of which is a design proposal. First, the road network is
automatically partitioned into coherent subregions using graph-based
community detection and density clustering (Louvain, Leiden, DBSCAN,
and four hybrid combinations), removing the need for manual,
expert-driven region design; this component is implemented and its
effect on the resulting region structure is measured across five
networks in Section~\ref{sec6}. Second, the Sub-Policy reward is
augmented with an explicit fairness penalty so that waiting-time
disparity, starvation, and Jain's Fairness Index are optimized
alongside local efficiency rather than left as an unmeasured
side-effect; this component is implemented (Section~\ref{sec4}) with
closed-loop training evaluation pending. Third, as a design proposal
rather than an implemented mechanism, the framework outlines how
single-vehicle emergency routing could be extended to a multi-EMV
setting in which several emergency vehicles simultaneously request
priority across overlapping parts of the network, together with a
mechanism for incident-robust training; neither is implemented in the
current codebase, and none of the five evaluation scenarios currently
contain emergency-vehicle demand. The remainder of this section
describes each of these three components in turn, with implementation
status stated explicitly at the start of each subsection.

% ---------------------------------------------------------------
% 3.1 CLUSTERING BASED ON LOUVAIN
% ---------------------------------------------------------------

\subsection{Clustering Based on Louvain}\label{subsec3-1}

\textit{Implementation status: implemented and evaluated
(Section~\ref{sec6}).}

In the proposed framework, manually dividing the road network into
precise subsections is impractical for real-world scenarios due to
the large number of roads and intersections involved. Therefore, the
road network is modeled as a weighted graph $G = (V, E)$, where
intersections are represented as nodes and road segments as edges.

To automatically identify suitable subsections, the Louvain
community detection algorithm (via the \texttt{python-louvain}
implementation) is applied to the graph. The algorithm
partitions the network into communities by maximizing
modularity, which measures the difference between the
observed connectivity within communities and the connectivity
expected in a corresponding random network:

\begin{equation}
Q = \frac{1}{2m}\sum_{i,j}\left(A_{ij} - \frac{k_i k_j}{2m}\right)
\delta(c_i, c_j),
\end{equation}

where $A_{ij}$ is the edge weight between nodes $i$ and $j$, $k_i$ and
$k_j$ are their weighted degrees, $m$ is the total edge weight, and
$\delta(c_i, c_j)$ indicates whether the nodes belong to the same
community.

A higher modularity indicates stronger community structure, allowing
the identified communities to serve as automatically generated
road-network subsections. Each subsection is assigned a
sub-policy, while a higher-level meta-policy
coordinates the behavior of these sub-policies across the complete
network. This provides a scalable alternative to manual decomposition
and enables hierarchical policy optimization over complex road
networks.

% ---------------------------------------------------------------
% 3.2 MULTIPLE EMV VEHICLE ROUTING
% ---------------------------------------------------------------

\subsection{Multiple EMV Vehicle Routing Across the City}\label{subsec3-2}

\textit{Implementation status: proposed design, not implemented. No
routing, conflict-arbitration, or priority-class code exists in the
current codebase, and none of the five evaluation scenarios (Section~
\ref{sec6}) include an emergency vehicle in their demand data. The
remainder of this subsection should be read as a design sketch for
future work, not as a description of an evaluated system.}

Prior emergency-vehicle (EMV) prioritization approaches in the
traffic-signal-control literature, including the single-agent
preemption strategies discussed in Section~\ref{sec2}, are almost
exclusively formulated around a single emergency vehicle traversing
the network at any given time. This assumption breaks down in dense
urban environments, where it is common for multiple ambulances, fire
trucks, and police vehicles to be simultaneously active, often
converging on overlapping segments of the road network from different
origins and with different destinations. Extending EMV prioritization
to a genuinely multi-vehicle setting therefore requires the framework
to reason jointly about several time-critical routes rather than
optimizing each one in isolation.

We propose to model the multi-EMV routing problem as a joint
optimization over a set of active emergency vehicles
$\mathcal{E} = \{e_1, e_2, \dots, e_K\}$, each associated with a
current position, a destination, and a priority class (for example,
life-threatening medical emergency versus structural fire response).
At each decision step, the framework must determine (i) a preferred
route for every active emergency vehicle through the signalized
network, and (ii) how to allocate green-phase priority at any
intersection that lies on more than one emergency vehicle's preferred
route. This second requirement is the central new difficulty relative
to single-EMV routing: when two emergency vehicles approach the same
intersection from conflicting directions within a short time window,
granting unconditional priority to both is physically impossible, and
the controller must instead arbitrate based on factors such as
estimated time-to-intersection, remaining distance to destination, and
priority class.

We envision extending the existing Meta-Policy/Sub-Policy hierarchy so
that each active emergency vehicle contributes an additional,
time-decaying bias term to the regional and intersection-level state
representations along its anticipated route, rather than introducing a
wholly separate routing module. Conceptually, this bias term would
increase in magnitude as an emergency vehicle approaches a given
intersection and decay once it has passed, so that the learned
Sub-Policy naturally favors clearing a path ahead of an approaching
emergency vehicle while reverting to its ordinary efficiency/fairness
trade-off once the vehicle has moved on. Coordination across
subregions, needed when an emergency vehicle's route crosses several
Louvain-derived subregions in sequence, would be handled by the
Meta-Policy, which is already responsible for issuing cross-regional
guidance.

Resolving conflicts between multiple simultaneously active emergency
vehicles is expected to require an explicit arbitration rule layered
on top of this bias mechanism --- for instance, prioritizing the
vehicle with the least estimated time-to-intersection, breaking ties
by priority class, and bounding the maximum delay that can be imposed
on a lower-priority emergency vehicle so that no single vehicle is
perpetually deprioritized in favor of others. The corresponding
evaluation methodology would extend the single-EMV response-time
metrics used elsewhere in this work (e.g., average and 95th-percentile
emergency waiting time) to the multi-vehicle case by additionally
reporting the spread of response-time improvements across all active
emergency vehicles, so that the framework can be shown to accelerate
every emergency vehicle's passage and not merely the one considered by
the controller to be of highest priority.

Finally, we note that a multi-EMV setting also changes how the
framework's fairness objective (Section~\ref{subsec3-3}) should be
interpreted: granting priority to emergency vehicles necessarily
imposes additional waiting time on ordinary traffic, and this
additional cost must be distributed across the network rather than
concentrated on whichever ordinary vehicles happen to be sharing a
road segment with an active emergency-vehicle corridor at a given
moment. This interaction between multi-EMV routing and the fairness
reward is left as an open design question to be resolved in the full
version of this section.

% ---------------------------------------------------------------
% 3.3 FAIRNESS AND ACCIDENTAL CONGESTION CONTROL
% ---------------------------------------------------------------

\subsection{Fairness and Accidental Congestion Control}\label{subsec3-3}

\textit{Implementation status: the fairness penalty described below
(variance, starvation threshold, Jain's Fairness Index) is implemented
and used in training (Section~\ref{sec4}); the envy term and the
accident/incident-injection mechanism described later in this
subsection are proposed designs, not implemented in the current
codebase.}

Traffic-signal controllers optimized purely for network-wide
efficiency metrics, such as average travel time or average delay, can
produce policies that are efficient in aggregate yet highly inequitable
at the level of individual approaches or vehicles. A controller may,
for instance, learn to sacrifice one chronically low-priority approach
in order to keep throughput high on the network's dominant corridors,
resulting in a small subset of vehicles experiencing persistent,
disproportionate waiting times even while the network-wide average
delay looks acceptable. The fairness component of the proposed
framework is intended to make this trade-off explicit and controllable
rather than an incidental by-product of reward design.

We plan to quantify fairness along several complementary axes: the
variance of waiting times across approaches at a given intersection,
the maximum (tail) waiting time experienced by any single approach,
and Jain's Fairness Index computed over per-approach or per-vehicle
service levels. Unlike a single scalar efficiency term, these measures
are each sensitive to a different failure mode --- variance captures
persistent imbalance between approaches, the tail-waiting-time penalty
specifically targets starvation of the worst-served approach, and
Jain's index provides a normalized, scale-independent summary of
overall equity that can be tracked and compared across scenarios and
partitioning methods. An ``envy'' term, comparing each intersection's
service level against that of its neighboring intersections, is
further intended to discourage locally greedy policies that improve
one intersection's performance strictly at a neighbor's expense ---
directly relevant to the Louvain-based regional decomposition of
Section~\ref{subsec3-1}, since poorly balanced subregions could
otherwise induce exactly this kind of neighbor-level inequity.

A second, largely orthogonal concern addressed by this subsection is
robustness to accidental congestion: collisions, breakdowns,
and other unplanned incidents that suddenly and sharply reduce the
effective capacity of one or more road segments. Conventional traffic
signal controllers are typically trained and evaluated exclusively
under nominal demand conditions, and so may respond poorly when a
lane or approach is unexpectedly blocked, continuing to allocate green
time as though the obstructed segment retained its usual capacity. We
intend to address this by exposing the controller to a representative
distribution of injected incidents during training --- varying in
location, duration, and severity --- so that the learned Sub-Policy
and Meta-Policy learn to detect the resulting shift in local traffic
statistics (e.g., a sudden rise in queue length or a collapse in mean
speed on an otherwise normal approach) and reallocate green time
accordingly, rather than requiring a separate, hand-engineered
incident-detection and override layer.

We additionally plan to evaluate whether accident-aware training
degrades the framework's fairness guarantees: an incident response
that simply redirects all available green time toward the obstructed
approach could, in principle, starve unrelated approaches elsewhere in
the network. The evaluation methodology for this subsection will
therefore report fairness metrics (variance of waiting time, tail
waiting time, Jain's index) separately for incident and incident-free
simulation intervals, to confirm that the controller's response to
accidental congestion does not come at the cost of the equity
objectives established earlier in this section. Together with the
multi-EMV routing extension of Section~\ref{subsec3-2}, this
completes the three methods --- scalable graph-based regional
decomposition, coordinated multi-emergency-vehicle routing, and
fairness-aware, incident-robust congestion control --- that constitute
the proposed FairLight framework.

% ---------------------------------------------------------------
% 3.4 HIERARCHICAL ARCHITECTURE OVERVIEW
% ---------------------------------------------------------------

\subsection{Hierarchical Architecture Overview}\label{subsec3-4}

\textit{Implementation status: unless otherwise noted, the
architecture described in Sections~\ref{subsec3-4}--\ref{subsec3-10}
(Meta-Policy, Sub-Policy, PPO training, and safety-constrained
execution) is implemented and verified against the training code;
closed-loop evaluation results are reported separately in
Section~\ref{sec6} and are pending at the time of writing.}

The proposed FairLight framework addresses the scalability and
coordination challenges of city-scale adaptive traffic signal control
through a hierarchical multi-agent reinforcement learning
architecture. A fully centralized controller becomes computationally
expensive as the number of intersections increases, whereas
independent intersection-level agents lack sufficient network-wide
context to coordinate traffic across connected corridors. To address
this trade-off, FairLight introduces an intermediate regional
abstraction between network-level coordination and local signal
control.

The framework consists of four principal components: (i)
automatic regional construction, (ii) a network-level Meta-Policy,
(iii) intersection-level Sub-Policies, and (iv) a safety-constrained
signal execution layer. Regional construction transforms the
signalized road network into coherent traffic regions, the Meta-Policy
extracts global spatial-temporal information and generates
network-level control targets, and the Sub-Policies translate this
global guidance into local phase decisions using intersection and
neighborhood information. The resulting hierarchy is trained jointly
using PPO with an adversarial goal-alignment mechanism.

% ---------------------------------------------------------------
% 3.5 TRAFFIC NETWORK REPRESENTATION AND AUTOMATIC REGIONALIZATION
% ---------------------------------------------------------------

\subsection{Traffic Network Representation and Automatic
Regionalization}\label{subsec3-5}

\textit{Implementation status: implemented and evaluated
(Section~\ref{sec6}). Region membership is computed once, offline,
before training begins, and is loaded from a static cluster file for
the entire training and evaluation run; only the edge weights feeding
Louvain and Leiden (below) are estimated from a traffic simulation
pass, not the community structure itself. This is a dynamic-weight,
static-membership design: it should not be described as continuously
adaptive regionalization.}

Manually decomposing a large urban road network into meaningful
control regions is impractical because the number and connectivity of
signalized intersections can vary substantially across networks.
FairLight therefore represents the signalized road network as a
weighted graph

\begin{equation}
G = (V, E),
\end{equation}

where each vertex $v_i \in V$ represents a signalized intersection and
each edge $e_{ij} \in E$ represents the connectivity between two
signalized intersections. Non-signalized junctions between traffic
lights are contracted during graph construction so that the resulting
graph directly represents the controllable signalized-intersection
topology.

The connectivity between intersections is characterized using a
traffic-dependent edge weight rather than purely static topology. For
an edge $e$, the implemented weighting function is

\begin{equation}
w_e = \alpha n_e + \beta \omega_e + \gamma \kappa_e n_e,
\end{equation}

where $n_e$ denotes the number of vehicles on the edge, $\omega_e$
represents accumulated waiting time, and

\begin{equation}
\kappa_e =
\operatorname{clip}\!\left(1 - \frac{v_e}{v_e^{\max}}, 0, 1\right)
\end{equation}

denotes the congestion ratio. The implementation uses $\alpha = 1.0$,
$\beta = 0.3$, and $\gamma = 2.0$. These weights are averaged over a
1000-step simulation window, allowing the regional structure to
reflect persistent traffic interactions rather than a transient
traffic state.

The resulting weighted graph is partitioned into communities using
graph community-detection methods. In particular, Louvain and Leiden
identify communities by maximizing graph modularity,

\begin{equation}
Q = \frac{1}{2m}\sum_{i,j}
\left(A_{ij} - \frac{k_i k_j}{2m}\right)
\delta(c_i, c_j),
\end{equation}

where $A_{ij}$ is the weighted adjacency matrix, $k_i$ and $k_j$ are
the weighted degrees of vertices $i$ and $j$, $m$ is the total edge
weight, and $\delta(c_i, c_j)$ indicates whether the two vertices
belong to the same community. Louvain is implemented via the
\texttt{python-louvain} package and Leiden via \texttt{leidenalg} with
the \texttt{ModularityVertexPartition} objective, both operating on the
traffic-weighted graph above; both post-process their raw output by
merging singleton or near-singleton communities into a neighboring
community before the partition is written out, so the reported region
counts in Section~\ref{sec6} already reflect this merge step.

DBSCAN is implemented as a separate, spatial clustering path: it
clusters intersections by their $(x,y)$ coordinates under Euclidean
distance, not by the traffic-weighted graph used for Louvain and
Leiden. Its \texttt{eps} and \texttt{min\_samples} parameters are
auto-tuned rather than fixed: \texttt{min\_samples} is set to the
network's node count divided by 10, clipped to $[3,8]$, and
\texttt{eps} is swept over a set of log-spaced candidates derived from
$k$-nearest-neighbor distances, selecting the candidate with the
lowest noise-point fraction whose resulting mean cluster size falls in
an $8$--$40$-node target band. The four hybrid methods combine the two
differently depending on composition order: \texttt{dbscan\_X} (for
community-detection method $X$) first applies spatial DBSCAN to form
fine-grained spatial cells and then runs $X$ on the resulting
cell-adjacency supernode graph, while \texttt{X\_dbscan} first runs
community detection $X$ on the traffic-weighted graph and then
subdivides each resulting community spatially using DBSCAN. Because
DBSCAN's clustering criterion is purely spatial, it does not directly
incorporate the traffic-dependent edge weight $w_e$ defined above; only
the Louvain, Leiden, and hybrid paths that invoke community detection
do so.

This regionalization produces a set of regions

\begin{equation}
\mathcal{R} = \{R_1, R_2, \ldots, R_M\},
\end{equation}

which forms the intermediate abstraction used by the hierarchical
controller. Each region provides aggregated traffic information to the
Meta-Policy, while individual intersections within the regions remain
independently controllable by the Sub-Policies.

% ---------------------------------------------------------------
% 3.6 HIERARCHICAL META-POLICY
% ---------------------------------------------------------------

\subsection{Hierarchical Meta-Policy}\label{subsec3-6}

The Meta-Policy provides network-level coordination by observing the
temporal evolution of traffic conditions across the generated regions.
For each region $R_k$, the regional state is represented as

\begin{equation}
s_k^t = \left[N_k^t, W_k^t, x_k, y_k\right] \in \mathbb{R}^{4},
\end{equation}

where $N_k^t$ denotes the aggregate stopped-vehicle count, $W_k^t$
denotes aggregate waiting time, and $(x_k, y_k)$ represents the
normalized spatial centroid of the region.

Because traffic conditions are inherently temporal, the Meta-Policy
maintains a sliding window of $T = 20$ regional observations. A
Transformer encoder processes the regional tokens jointly to capture
spatial dependencies between regions. A learnable global token is
introduced to obtain a network-level representation

\begin{equation}
F_g \in \mathbb{R}^{4}.
\end{equation}

The temporal representations of individual regions are subsequently
processed using an LSTM. The resulting hidden states are concatenated
across the $M$ regions and transformed through an MLP to obtain a
global representation $G \in \mathbb{R}^{16}$. The Meta-Policy then
generates two network-level control targets,

\begin{equation}
(G_w, G_q) = \operatorname{Linear}_{16\rightarrow2}(G),
\end{equation}

where $G_w$ represents the desired global waiting-time level and
$G_q$ represents the desired global queue-length level.

Thus, rather than directly selecting signal phases throughout the
network, the Meta-Policy establishes a global traffic-control
objective that can be pursued by the lower-level intersection agents.
This separation allows long-range traffic coordination to be handled
at a compressed regional level while preserving local responsiveness
at individual intersections.

% ---------------------------------------------------------------
% 3.7 INTERSECTION-LEVEL SUB-POLICY
% ---------------------------------------------------------------

\subsection{Intersection-Level Sub-Policy}\label{subsec3-7}

\textit{Implementation status: implemented. The neighborhood mechanism
described below is a fixed-capacity, top-4 attention layer, not an
arbitrary-degree mechanism; see the clarification after
Eq.~\eqref{eq:gac-concat}.}

The Sub-Policy converts global guidance into local signal-control
actions. At each intersection $i$, the local observation is
represented by

\begin{equation}
o_i^t = \left[v_i, q_i, o_i, f_i, s_i, w_i, v_i^{avg}, p_i, c_i,
d_i\right] \in \mathbb{R}^{10},
\end{equation}

where the components capture local traffic volume, queue length,
occupancy, flow, stops, waiting time, speed, pressure, congestion, and
delay. The observations are normalized using an exponential moving
average and clipped to a bounded range before being supplied to the
policy network.

A local multilayer perceptron first encodes the intersection state.
Since traffic conditions at one intersection are influenced by
neighboring intersections, the encoded representation is augmented
using a Graph Attention Concat (GAC) mechanism over the
intersection-adjacency graph. For a neighboring pair $(i, j)$, the
attention coefficient is computed as

\begin{equation}
e_{ij} = \operatorname{LeakyReLU}\!\left(a^{T}[h_i \Vert h_j]\right),
\qquad
\alpha_{ij} = \operatorname{softmax}_j(e_{ij}),
\end{equation}

and the resulting neighborhood representation is constructed through
concatenation,

\begin{equation}
z_i = h_i \Vert \alpha_{i1}h_{j_1} \Vert \cdots \Vert
\alpha_{i4}h_{j_4}.
\label{eq:gac-concat}
\end{equation}

The neighbor set $\{j_1,\ldots,j_4\}$ in Eq.~\eqref{eq:gac-concat} is a
fixed-capacity slot of exactly four, not a variable-length
representation of an intersection's true adjacency degree: when an
intersection has more than four signalized neighbors, the remainder are
truncated from the representation; when it has fewer than four, the
unused slots are zero-padded. The regional partitioning of
Section~\ref{subsec3-5} determines \emph{which} intersections are
grouped into a region, but the attention mechanism itself operates on
this fixed four-neighbor window regardless of an intersection's actual
degree in the signalized road network. Characterizing this mechanism as
adapting to arbitrary network degree would overstate what is
implemented; its correct characterization is a capped, zero-padded
neighborhood of size four.

The neighborhood representation is then combined with the global
Meta-Policy embedding:

\begin{equation}
s_i^t = \left[z_i \Vert F_g\right].
\end{equation}

This representation simultaneously provides the Sub-Policy with
local traffic information, neighboring traffic context, and
network-level coordination guidance.

A shared actor-critic network processes this representation and
produces a probability distribution over signal phases. Because
intersections may have different numbers of valid signal phases,
invalid action indices are masked before the final softmax operation.
Parameter sharing across intersections allows the same learned policy
structure to be deployed across heterogeneous signalized intersections
while retaining local observations and graph-based neighborhood
information.

% ---------------------------------------------------------------
% 3.8 JOINT HIERARCHICAL OPTIMIZATION
% ---------------------------------------------------------------

\subsection{Joint Hierarchical Optimization}\label{subsec3-8}

The Meta-Policy and Sub-Policies are optimized jointly. The Sub-Policy
is trained using Proximal Policy Optimization (PPO) with Generalized
Advantage Estimation (GAE). The temporal-difference residual is

\begin{equation}
\delta_t = r_t + \gamma V(s_{t+1})(1 - d_t) - V(s_t),
\end{equation}

and the generalized advantage estimate is

\begin{equation}
A_t = \delta_t + \gamma\lambda(1 - d_t)A_{t+1},
\end{equation}

with $\gamma = 0.99$ and $\lambda = 0.95$. The PPO policy objective
uses the clipped probability ratio

\begin{equation}
L_{\text{policy}} = -\min\left(r_tA_t,
\operatorname{clip}(r_t, 1-\epsilon, 1+\epsilon)A_t\right),
\end{equation}

where $\epsilon = 0.2$.

The hierarchical relationship is enforced through an additional
goal-alignment term. Let $W_{\mathrm{global}}$ and $Q_{\mathrm{global}}$
denote the actual network-wide waiting time and queue length. The
Meta-Policy is trained using

\begin{equation}
L_{\mathrm{Meta}} = \operatorname{MSE}\!\left((G_w, G_q),
(W_{\mathrm{global}}, Q_{\mathrm{global}})\right) + \eta_1 r_g,
\end{equation}

while the Sub-Policy receives an alignment objective encouraging the
realized network state to follow the Meta-Policy's targets:

\begin{equation}
L_{\mathrm{Sub}} = L_{\mathrm{AC}} + \eta_2 \left[
\beta_q(Q_{\mathrm{global}} - G_q) +
\beta_w(W_{\mathrm{global}} - G_w)\right].
\end{equation}

Here, $\eta_1 = \eta_2 = 0.1$ and $\beta_q = \beta_w = 0.5$. The
alignment gradient is detached from the Meta-Policy in the Sub-Policy
loss, preventing the lower-level policy from directly modifying the
higher-level representation through this term.

The local reward combines intersection-level traffic efficiency with
the global goal-alignment reward:

\begin{equation}
r_i^t = -\left(q_i^t + w_i^t + d_i^t + p_i^t - s_i^t\right),
\end{equation}

\begin{equation}
r_g^t = -\left[\beta_q(W_{\mathrm{global}} - G_q) +
\beta_w(Q_{\mathrm{global}} - G_w)\right],
\end{equation}

and

\begin{equation}
r^t = r_i^t + r_g^t.
\end{equation}

This formulation ensures that each intersection is evaluated according
to its own local traffic conditions while simultaneously receiving a
network-level coordination signal.

% ---------------------------------------------------------------
% 3.9 SAFETY-CONSTRAINED SIGNAL EXECUTION
% ---------------------------------------------------------------

\subsection{Safety-Constrained Signal Execution}\label{subsec3-9}

The learned policy does not directly bypass traffic-signal safety
constraints. Every requested phase transition is passed through a
deterministic execution layer that enforces a minimum green interval
followed by a yellow transition before the new green phase is
activated.

Specifically, a requested phase change is permitted only after the
current phase has satisfied the minimum-green requirement. The
controller then enters a yellow transition for the prescribed duration
before committing the requested green phase. In the current
implementation, the minimum green time is 5 simulated seconds and the
yellow transition lasts 3 seconds.

This separation between learned decision-making and
deterministic safety enforcement prevents the reinforcement-learning
agent from producing physically invalid or unrealistically rapid
signal transitions.

% ---------------------------------------------------------------
% 3.10 END-TO-END CONTROL PROCESS
% ---------------------------------------------------------------

\subsection{End-to-End Control Process}\label{subsec3-10}

At every reinforcement-learning decision step, the framework first
aggregates the current traffic state for each automatically generated
region and updates the temporal regional-state buffer. The Meta-Policy
then processes this temporal representation to obtain the global
embedding and network-level queue and waiting-time targets.
Simultaneously, each intersection obtains its normalized local
observation, incorporates information from neighboring intersections
through graph attention, and fuses this representation with the global
Meta-Policy embedding.

The resulting Sub-Policy actions are passed through the
safety-constrained signal execution layer before being applied to the
traffic simulator. The environment then advances the traffic
simulation, computes local and global rewards, and stores the
resulting transition for policy optimization. Meta-Policy updates are
performed periodically, while Sub-Policy PPO updates are performed
after the rollout buffer reaches the specified horizon.

The complete architecture can therefore be expressed as the following
hierarchical mapping:

\begin{equation}
\begin{aligned}
G \rightarrow \mathcal{R} \rightarrow \text{Meta-Policy} \rightarrow
F_g, (G_w, G_q) \\
\rightarrow \text{Sub-Policies} \rightarrow \text{Signal Actions}
\rightarrow \text{Traffic Environment}.
\end{aligned}
\end{equation}

This design separates regional coordination, temporal network
reasoning, local spatial interaction, and physical signal execution
into distinct layers. Consequently, the framework can automatically
adapt its regional abstraction to different road networks while
retaining a shared local control mechanism and a global coordination
signal.


% ===============================================================
% MULTI-OBJECTIVE REWARD FUNCTION
% ===============================================================

\section{Multi-Objective Reward Function Design}\label{sec4}

An effective traffic signal controller must optimize multiple
objectives simultaneously rather than maximizing traffic throughput
alone. Conventional reinforcement learning approaches primarily
optimize queue lengths or average waiting times, often producing
policies that systematically disadvantage certain intersections and
fail to prioritize emergency vehicles. This section first states the
design target for a complete multi-objective reward, then states
precisely which part of that design target is implemented in the
current training code, to avoid the two being conflated.

\textbf{Design target.} The intended overall reward at each decision
step is

\begin{equation}
R_{\text{total}}
=
w_1 R_{\text{eff}}
+
w_2 R_{\text{fair}}
+
w_3 R_{\text{emergency}},
\end{equation}

with $w_1 + w_2 + w_3 = 1$ and $w_1 \ge w_2 \ge w_3 \ge 0$, assigning
$R_{\text{eff}}$ the largest weight because it governs global
traffic-congestion control across the whole network rather than
emergency-vehicle service alone, followed by fairness and then
emergency responsiveness.

\textbf{As implemented.} The weighted combination above, with explicit
$w_1,w_2,w_3$ coefficients, is \emph{not} present in the training code:
no such weight variables exist, and $R_{\text{emergency}}$ is not added
to the reward the Sub-Policy is trained on. What \emph{is} implemented
and used for training is the local reward of
Section~\ref{subsec3-8}, $r_i^t = -(q_i^t+w_i^t+d_i^t+p_i^t-s_i^t)$,
an absolute (not temporal-difference) combination of queue length,
waiting time, delay, pressure, and speed, plus the fairness penalty
$R_{\mathrm{fair},j}$ defined later in this section, plus the
goal-alignment term $r_g^t$. The discussion below presents
$R_{\text{eff}}$ and $R_{\text{emergency}}$ as proposed refinements to
be reconciled with $r_i^t$ in future work, not as a description of the
currently trained objective.

\subsection{Traffic Efficiency Reward (proposed refinement)}

\textit{Implementation status: not implemented as stated below; the
efficiency term actually used in training is $r_i^t$ from
Section~\ref{subsec3-8}.} The following formulation is proposed as a
denser-signal refinement of $r_i^t$, encouraging continuous improvement
in network performance by rewarding reductions in congestion rather
than rewarding only low congestion states:

\begin{equation}
R_{\text{eff}}
=
-\beta_1 \Delta Q
-\beta_2 \Delta W
+
L_{\text{avg}},
\end{equation}

where $\Delta Q$ denotes the change in total queue length between
successive decision steps, $\Delta W$ represents the change in
cumulative vehicle waiting time, and $L_{\text{avg}}$ is the average
local throughput reward aggregated across all intersections. The
coefficients $\beta_1$ and $\beta_2$ determine the relative importance
assigned to queue reduction and waiting-time reduction.

Unlike conventional reward formulations based solely on absolute queue
lengths, the use of temporal differences would in principle provide
denser learning signals by rewarding measurable improvements between
consecutive time steps; whether this translates into faster convergence
in practice has not been tested, since this term is not yet wired into
the training loop.

\subsection{Fairness Reward}

\textit{Implementation status: implemented. The coefficients and
thresholds below are literal values in the training code
(\texttt{environment.py}), verified against the implementation.}

Conventional traffic-signal controllers commonly optimize aggregate measures such as queue
length, average delay, and throughput. These objectives can improve network efficiency but do
not guarantee equitable service to individual vehicles: a movement repeatedly deprioritized by
the controller may experience prolonged waiting even when network-level metrics improve.
FairLight therefore treats fairness as a complement to traffic efficiency and emergency-vehicle
responsiveness, consistent with the multi-objective reward defined above.

\paragraph{Vehicle waiting-time monitoring.}
For each vehicle $v$ and intersection $j$, let $w_{v,j}(t)$ denote its accumulated waiting time
at time $t$. The timer increases by one simulation step whenever the vehicle is stopped (speed
below $0.1$~m/s) on an incoming lane of $j$, is kept while the vehicle creeps forward in the
queue, and is discarded when the vehicle is served and leaves the incoming lanes of $j$; a new
waiting episode then starts at the next intersection. The same definition is used in training
and evaluation. A vehicle is flagged as a starvation risk when
\begin{equation}
w_{v,j}(t)\ge W_{\max},\qquad W_{\max}=90\ \text{s}.
\end{equation}
Reaching $W_{\max}$ flags the vehicle and increases its consideration in subsequent phase
selection: the requested phase of the Sub-Policy is replaced by a feasible phase that gives a
green to the lane of the longest-waiting flagged vehicle (the current phase is kept if it
already does). The 90-second threshold is a design threshold and not a universal guarantee
that every vehicle is served within 90 seconds: demand, conflicting movements, incidents,
clearance intervals, and physical feasibility may make the threshold temporarily or
persistently unattainable. The override only changes the \emph{requested} phase; the controller
never bypasses the minimum-green requirement (5 steps) or the 3~s yellow clearance of
Section~\ref{subsec3-9}, and no override is issued during a yellow transition.

To align this safeguard with the reward, define the local maximum waiting time at intersection
$j$ as $W_j(t)=\max_{v\in\mathcal{V}_j(t)} w_{v,j}(t)$, where $\mathcal{V}_j(t)$ is the set of
tracked vehicles currently waiting at that intersection. The threshold penalty is then
$\max(0,\,W_j(t)-W_{\max})$. The maximum is taken locally per intersection, not over the whole
episode or network.

\paragraph{Intersection-level fairness using Jain's index.}
At each intersection, Jain's Fairness Index (JFI) summarizes how evenly a set of non-negative
service scores is distributed. For $n$ included vehicles with scores $x_i\ge 0$,
\begin{equation}
J=\frac{\left(\sum_{i=1}^{n}x_i\right)^2}{n\sum_{i=1}^{n}x_i^{2}}.
\end{equation}
For non-zero scores, $J\in[1/n,1]$, and a value closer to $1$ indicates a more even
distribution. If all scores are zero the denominator vanishes; we adopt the convention
$J=1$ (no disparity) and, when no vehicle is waiting at all, set every fairness term to zero.

The choice of $x_i$ matters. Raw waiting time should not be used as a service score, because a
higher waiting time indicates \emph{worse} service, and equal but very high waiting times still
yield $J=1$. We therefore use the bounded, non-negative score that is inversely related to
waiting time,
\begin{equation}
x_i=\frac{1}{1+w_{i,j}/W_{\mathrm{ref}}},
\end{equation}
where $W_{\mathrm{ref}}=30$~s is a fixed reference time chosen before training. The score, the
population (all vehicles waiting on the incoming lanes of $j$), the sampling interval (every
simulation step), and the handling of vehicles not yet served are identical across methods and
experiments. JFI measures equality, not absolute performance or starvation by itself, so it is
reported alongside the mean, maximum, and 95th-percentile waiting time, queue length,
throughput, and emergency-vehicle response time, and it is computed locally per intersection
before any city-wide aggregation so that local disparities are not concealed by a network
average.

\paragraph{Fairness reward.}
Let $W_j$ be the maximum waiting time at intersection $j$ as defined above, let
$\sigma^2_{\hat w,j}$ be the variance of the normalized waiting times
$\hat w_i=w_{i,j}/W_{\mathrm{ref}}$ over the local population, and let $J_j$ be the JFI of the
scores $x_i$. The dimensionless local fairness penalty is
\begin{equation}
P_{\mathrm{fair},j}=\lambda_{\mathrm{var}}\,\mathrm{Var}(\hat w)
+\lambda_{\mathrm{starve}}\max\!\Big(0,\frac{W_j-W_{\max}}{W_{\max}}\Big)
+\lambda_{J}\,(1-J_j),
\end{equation}
and the fairness reward is $R_{\mathrm{fair},j}=-P_{\mathrm{fair},j}$, which is added to the
local reward $r_j^t$ of Section~\ref{subsec3-8}. Normalizing waiting times by $W_{\mathrm{ref}}$
and the starvation excess by $W_{\max}$ is important because waiting-time variance has
squared-time units whereas the starvation and JFI terms are dimensionless. The coefficients
$\lambda_{\mathrm{var}},\lambda_{\mathrm{starve}},\lambda_J\ge 0$ are hyperparameters; larger
disparity, threshold excess, or inequality reduces the reward. Their initial values
$(\lambda_{\mathrm{var}},\lambda_{\mathrm{starve}},\lambda_J)=(1,10,10)$ were chosen so that the
fairness terms are not negligible next to the queue and waiting terms of $r_j^t$; they are
design choices, not empirically established optimal values. The earlier coefficients $0.3$, $0.25$ and $0.15$ applied to
quantities with incompatible units and are superseded by this formulation. Envy and
pedestrian terms are not part of the implemented fairness reward.

\subsection{Emergency Vehicle Reward (proposed, not wired into training)}

\textit{Implementation status: not implemented. The environment code
includes a read-only helper that computes cumulative emergency-vehicle
waiting time for monitoring purposes, but this quantity is never added
to the reward the policy is trained on, and no scenario's current
demand data contains an emergency vehicle for it to measure. The
formulation below is therefore a proposed reward term, not a
description of current training behavior.}

Emergency vehicles would require substantially higher priority than
regular traffic because delays directly affect emergency response
times. The proposed reward component is

\begin{equation}
R_{\text{emergency}}
=
-W_{\text{emergency}},
\end{equation}

where $W_{\text{emergency}}$ is the cumulative waiting time
experienced by emergency vehicles. This formulation would directly
penalize any delay encountered by ambulances, fire trucks, or police
vehicles, encouraging the controller to allocate green phases
immediately upon emergency vehicle detection, once (i) emergency
vehicles are added to the simulated demand and (ii) this term is
connected to the training loop.

Finally, the combined reward is normalized using an exponential
moving average before being supplied to the reinforcement learning
algorithm. Reward normalization reduces reward variance, improves
numerical stability during training, and accelerates policy
convergence.

\subsection{Evaluation Protocol}

Fairness is evaluated jointly with efficiency and safety. At minimum we report the mean,
maximum, and 95th-percentile vehicle waiting time; the proportion of vehicles exceeding
90~seconds; intersection-level JFI; queue length; throughput; and emergency response time. The
controller is compared against the same controller without fairness terms, under identical
traffic demand, random seeds where applicable, signal constraints, and simulation horizons, and
multiple runs are reported with variability estimates. A reduction in average delay alone is
insufficient evidence of improved fairness, and a higher JFI alone is insufficient unless
absolute delays remain acceptable. Emergency pre-emption intentionally gives some vehicles
higher urgency, so results distinguish ordinary-vehicle fairness from emergency-vehicle
responsiveness. At present, the training and evaluation code does not set a random seed for
either the simulator or the learning algorithm (no seeding call exists in the training
entry point), so repeated runs under this protocol are not currently reproducible run-to-run;
adding explicit seed control is a prerequisite for the multi-run variability estimates this
protocol calls for, and is tracked as a limitation in Section~\ref{sec7}.


% ===============================================================
% EXPERIMENTATION
% ===============================================================

\section{Experimentation}\label{sec5}

\subsection{Simulation Environment}\label{subsec5-1}

The proposed framework is evaluated in the Simulation of Urban MObility
(SUMO) environment via the TraCI interface. SUMO is chosen because its
microscopic, lane-level vehicle dynamics and configurable signal
infrastructure allow the simulated network to closely approximate the
congestion and coordination patterns of a real signalized road network,
providing a controlled setting in which the proposed regional
partitioning, fairness, and emergency-routing mechanisms can be
exercised repeatedly under known traffic demand. A simulation study,
however, remains an approximation: real-world deployment would need to
contend with sensor noise, imperfect vehicle detection, driver
non-compliance, and infrastructure constraints that are not fully
captured by SUMO's vehicle models. Validating the framework on a real
traffic signal system is therefore treated as a necessary step beyond
the present simulation-based evaluation, and not as a substitute for
it.

\subsection{Base Model and Benchmark Methods}\label{subsec5-2}

The base model used throughout this evaluation is the Hierarchical
Region-Aware Multi-Agent Reinforcement Learning controller with the
fairness reward of Section~\ref{subsec3-3} enabled (Meta-Policy +
Sub-Policy, trained with the multi-objective reward of
Section~\ref{sec4}). This base model is compared against two
traditional, non-learning traffic signal control strategies commonly
used as references in the traffic-signal-control literature:

\begin{itemize}
\item \textit{Fixed Time}: a cyclical, predefined signal-phase schedule
with a fixed green/yellow split and no responsiveness to real-time
traffic conditions.
\item \textit{Max Pressure}: a reactive, non-learning controller that
selects the signal phase maximizing the pressure difference between
incoming and outgoing lanes at each intersection~\cite{maxpressure}.
\end{itemize}

The city-wide multi-EMV routing comparison described in some of this
framework's design documents is \emph{not} currently executable: the
routing and arbitration mechanism of Section~\ref{subsec3-2} is
unimplemented, and none of the five scenarios below include an
emergency vehicle in their demand data, so there is no EMV traffic for
an ``with vs.\ without city-wide routing'' comparison to act on. Where
such a comparison appears in Section~\ref{sec6} it is explicitly marked
as blocked on these two prerequisites, not merely pending compute time.

\subsection{Traffic Network Scenarios and Experimental Configuration}\label{subsec5-2b}

The framework is configured for five SUMO networks, selected to vary
network scale, topology regularity, and provenance (real-world OSM
extracts versus synthetic grids). Table~\ref{tab:scenarios} reports
the verified configuration of each, obtained directly from the
\texttt{.net.xml} and demand files shipped in the repository; entries
are marked ``Not reported'' where the corresponding configuration file
does not specify a value, rather than inferring one.

\begin{table}[t]
\centering
\caption{Traffic network scenarios used for regionalization evaluation
(Section~\ref{sec6}). Intersection counts are signalized (traffic-light)
junctions in the contracted intersection graph of
Section~\ref{subsec3-5}; vehicle counts are the configured demand for
the default route file. None of the five scenarios currently include an
emergency vehicle in their demand data.}
\label{tab:scenarios}
\begin{tabular}{lccccc}
\toprule
Scenario & Intersections & Edges & Vehicles & Sim.\ window [s] & Network type \\
\midrule
Manhattan     & 472 & 6184 & 1800 & Not reported   & OSM, real-world \\
Cologne8      & 8   & 590  & 2046 & 25200--28800 (3600) & OSM, real-world \\
Ingolstadt21  & 21  & 2415 & 4283 & 57600--61200 (3600) & OSM, real-world \\
Arterial4x4   & 16  & 320  & 2484 & 0--3600 (3600) & Synthetic, arterial grid \\
Grid4x4       & 16  & 464  & 1473 & 0--3600 (3600) & Synthetic, grid \\
\bottomrule
\end{tabular}
\end{table}

Arterial4x4 and Grid4x4 each ship 1400 pre-generated, numbered demand
variants in addition to the default (\texttt{\_1}) file reported above;
these appear intended as seed/demand variants for repeated runs but are
not yet invoked by any automated sweep in the training or evaluation
code, so the single-run vehicle counts in Table~\ref{tab:scenarios}
reflect only the default configuration. EMV origin, destination, and
priority-class fields, and the number of regional communities generated
per scenario, are reported in Section~\ref{sec6} (the latter) or marked
not applicable (the former, since no scenario currently defines an EMV).

\subsection{Neighborhood Construction}\label{subsec5-3}

The Sub-Policy's Graph Attention Concat mechanism (Section~\ref{subsec3-7},
Eq.~\eqref{eq:gac-concat}) operates on a fixed-capacity window of the
four highest-relevance signalized neighbors per intersection, with
truncation above four neighbors and zero-padding below four. This is
not a variable-degree or topology-adaptive mechanism: it is capped
regardless of an intersection's true degree in the signalized road
network. Whether any of the five scenarios in
Table~\ref{tab:scenarios} contain signalized intersections with more
than four neighbors --- and therefore lose information to truncation
--- has not been audited at the time of writing; this is listed as an
open question in Section~\ref{sec7}.

\subsection{Evaluation Metrics}\label{subsec5-4}

Four categories of metric are used, following the structure of
Section~\ref{sec6}. Only metrics collectible from the existing TraCI
instrumentation or from the static clustering output are reported as
measured; the remainder require a completed training/evaluation run
and are reported as pending.

\textit{Traffic efficiency}: total average trip time (vehicles
completing their trip within the simulation horizon), total waiting
time, and total average waiting time, consistent with
Section~\ref{subsec3-8}'s reward terms.

\textit{Emergency responsiveness}: mean, median, and 95th-percentile
EMV response time, and route-completion rate --- all currently
inapplicable, since no scenario contains an EMV.

\textit{Ordinary-vehicle fairness}: Jain's Fairness Index, mean and
95th-percentile waiting time, the proportion of vehicles exceeding the
$W_{\max}=90$~s threshold, and per-intersection waiting-time variance,
as defined in Section~\ref{sec4}.

\textit{Scalability and regionalization}: number of detected regions,
region-size distribution, and modularity per partitioning method. These
are computed directly from the static clustering output and do not
require a trained policy; they are reported as measured results in
Section~\ref{sec6}.


% ===============================================================
% RESULTS AND DISCUSSION
% ===============================================================

\section{Results and Discussion}\label{sec6}

\subsection{Regionalization Comparison (measured)}\label{subsec6-0}

Unlike the tables later in this section, Table~\ref{tab:regionalization}
reports measured output of the implemented clustering pipeline
(Section~\ref{subsec3-5}) applied to all five scenarios of
Table~\ref{tab:scenarios}: no trained policy is required to produce
these numbers, so they are reported as actual results rather than
placeholders.

\begin{table}[t]
\centering
\caption{Regions produced per scenario and method: region count $n$,
mean region size (intersections/region), and modularity $Q$ where
applicable (DBSCAN does not optimize modularity, since it clusters on
spatial coordinates rather than the traffic-weighted graph;
Section~\ref{subsec3-5}). Values are the current repository's cluster
output.}
\label{tab:regionalization}
\begin{tabular}{lccc}
\toprule
Scenario & DBSCAN ($n$, mean size) & Louvain ($n$, mean size, $Q$) & Leiden ($n$, mean size, $Q$) \\
\midrule
Manhattan    & 8, 59.0 & 9, 52.4, 0.620  & 70, 6.7, 0.577 \\
Cologne8     & 1, 8.0  & 1, 8.0, 0.000   & 2, 4.0, 0.012 \\
Ingolstadt21 & 2, 10.5 & 1, 21.0, 0.005  & 2, 10.5, 0.052 \\
Arterial4x4  & 1, 16.0 & 2, 8.0, 0.414   & 5, 3.2, 0.153 \\
Grid4x4      & 1, 16.0 & 1, 16.0, 0.375  & 5, 3.2, 0.288 \\
\bottomrule
\end{tabular}
\end{table}

These results answer RQ1 (Section~\ref{sec1}): the choice of
regionalization method materially changes the structure handed to the
Meta-Policy, and the effect is network-dependent rather than uniform.
On Manhattan --- the largest, most topologically rich network ---
Leiden produces far more, far smaller regions than Louvain (70 vs.\ 9),
with 34 of Leiden's 70 regions being singletons despite the
singleton-merging post-processing of Section~\ref{subsec3-5}; both
achieve comparable, reasonably high modularity ($0.577$--$0.620$). On
the three smaller or sparser networks (Cologne8, Ingolstadt21), Louvain
and DBSCAN both frequently collapse to a single region covering the
entire network, and modularity is close to zero (Cologne8 Louvain:
$Q{=}0$) --- indicating these networks lack strong community structure
at this scale, so ``automatic regionalization'' degenerates to no
regionalization at all for them under Louvain. DBSCAN, which clusters
purely on intersection coordinates rather than traffic load, produces
coarser and more size-uniform partitions across all five networks
than either community-detection method, consistent with it optimizing
a different (spatial) criterion rather than a worse one. These results
indicate that the choice of regionalization method should be
validated per-network rather than fixed globally, and that a method
producing many small regions on one network (Leiden on Manhattan) can
produce almost none on another (Leiden on Cologne8), which is
important context for interpreting any downstream control-performance
difference once closed-loop results are available.

\subsection{Effect of City-Wide EMV Routing (blocked, not merely pending)}\label{subsec6-1}

Table~\ref{tab:emv-routing} is the comparison this framework would
support between the base model operated with city-wide multi-EMV
routing (Section~\ref{subsec3-2}) and the same base model without it.
Unlike Table~\ref{tab:baseline-comparison}, this comparison is
currently blocked on two unmet prerequisites rather than simply awaiting
compute time: the routing and arbitration mechanism itself is
unimplemented, and no scenario's demand data contains an emergency
vehicle for the comparison to act on. The entries below are
placeholders recorded to make the intended comparison explicit; they
must not be read as measured results, and should not be populated until
both prerequisites are resolved.

\begin{table}[t]
\centering
\caption{Base model with vs.\ without city-wide EMV routing --- placeholder structure; blocked on unimplemented routing code and the absence of EMV demand in all five scenarios (Table~\ref{tab:scenarios}).}
\label{tab:emv-routing}
\small
\begin{tabular}{lccc}
\toprule
Configuration & Trip Time [s] & Waiting Time [s] & Avg.\ Waiting Time [s] \\
\midrule
Without city-wide EMV routing & 0.00 & 0.00 & 0.00 \\
With city-wide EMV routing    & 0.00 & 0.00 & 0.00 \\
\bottomrule
\end{tabular}
\end{table}

\subsection{Comparison with Baseline Traffic Control Methods (pending)}\label{subsec6-2}

Table~\ref{tab:baseline-comparison} reports the three evaluation
metrics of Section~\ref{subsec5-4} for Fixed Time, Max Pressure, and
the proposed base model (Algorithm). Unlike
Section~\ref{subsec6-1}, this comparison is not blocked by missing
functionality --- Fixed Time, Max Pressure, and the PPO-trained base
model are all implemented --- but the corresponding simulation runs
have not yet been executed under the current (PPO-based) training code
at the time of writing, and the codebase does not currently set a
random seed (Section~\ref{sec4}), so a single run would not meet this
paper's own evaluation protocol. All entries below are placeholders
pending a seeded, multi-run evaluation campaign.

\begin{table}[t]
\centering
\caption{Comparison of traffic signal control methods --- placeholder values pending a completed, seeded, multi-run evaluation campaign.}
\label{tab:baseline-comparison}
\small
\begin{tabular}{lccc}
\toprule
Method & Trip Time [s] & Waiting Time [s] & Avg.\ Waiting Time [s] \\
\midrule
Fixed       & 0.00 & 0.00 & 0.00 \\
Max Pressure & 0.00 & 0.00 & 0.00 \\
Algorithm   & 0.00 & 0.00 & 0.00 \\
\bottomrule
\end{tabular}
\end{table}

\subsection{Discussion and Limitations}\label{subsec6-3}

The one populated result in this section, Table~\ref{tab:regionalization},
already motivates a methodological caution: because regionalization
behavior is network-dependent (Section~\ref{subsec6-0}), any future
closed-loop comparison between regionalization methods must be
reported per-scenario, not pooled, since a method that is clearly
better on Manhattan (e.g.\ Leiden's higher modularity) may be
indistinguishable from the others on Cologne8 or Ingolstadt21 where all
methods nearly collapse to one region.

The remaining, unpopulated tables reflect genuine gaps rather than
omissions: Table~\ref{tab:emv-routing} cannot be populated until the
multi-EMV mechanism of Section~\ref{subsec3-2} is implemented and at
least one scenario's demand file includes an emergency vehicle;
Table~\ref{tab:baseline-comparison} requires a seeded, multi-run
training and evaluation campaign that has not yet been executed under
the current PPO-based code. No performance, efficiency, or
fairness-improvement claim is made on the basis of either table's
placeholder entries. Additional limitations established in this audit
--- the fixed four-neighbor attention cap (Section~\ref{subsec5-3}),
the gap between the proposed weighted reward design and the reward
actually trained on (Section~\ref{sec4}), and the absence of random
seeding --- are consolidated in Section~\ref{sec7}.


% ===============================================================
% CONCLUSION
% ===============================================================

\section{Conclusion}\label{sec7}

This paper presented FairLight, a hierarchical multi-agent
reinforcement learning framework for traffic signal control. Three
things are actually implemented and verified against the training
code: (i) automatic, graph-based regional partitioning via Louvain,
Leiden, DBSCAN, and four hybrid combinations, measured across five SUMO
networks in Section~\ref{sec6}; (ii) a hierarchical Meta-Policy/
Sub-Policy architecture trained with PPO and GAE, with a fixed,
zero-padded four-neighbor graph-attention mechanism at the Sub-Policy
level; and (iii) an explicit ordinary-traffic fairness penalty
(waiting-time variance, a 90-second starvation threshold, and Jain's
Fairness Index) added to the Sub-Policy's local reward. What the
experiments in this paper establish is narrower than the full research
programme: Table~\ref{tab:regionalization} shows that regionalization
method choice produces substantially different, network-dependent
region structures (from 1 region with near-zero modularity on
Cologne8/Ingolstadt21 under Louvain, to 70 regions on Manhattan under
Leiden); closed-loop performance comparisons against Fixed-Time and
Max-Pressure baselines, and any comparison involving city-wide EMV
coordination, remain pending or blocked, respectively, and are not
claimed here.

Several limitations currently prevent stronger conclusions. The
weighted multi-objective reward of Section~\ref{sec4}
($w_1 R_{\text{eff}} + w_2 R_{\text{fair}} + w_3 R_{\text{emergency}}$)
is a design target, not the reward the training code optimizes, which
is instead the absolute local reward of Section~\ref{subsec3-8} plus
the implemented fairness penalty; this gap should be closed before any
claim is made about efficiency/fairness/emergency trade-offs under the
stated weights. The multi-EMV routing design of
Section~\ref{subsec3-2} and the incident-robustness mechanism of
Section~\ref{subsec3-3} are proposed but unimplemented, and no
evaluation scenario currently contains an emergency vehicle, so no
EMV-related claim --- response time, fairness cost of preemption, or
otherwise --- is currently testable in this codebase. The fixed
four-neighbor attention cap (Section~\ref{subsec5-3}) may discard
information at intersections with higher signalized degree; whether
this occurs in the five evaluated networks has not been audited. The
training and evaluation code does not set a random seed, so the
multi-run variability estimates this paper's own evaluation protocol
calls for (Section~\ref{sec4}) cannot yet be produced. Finally, an
earlier, structurally different (pre-PPO) version of this codebase was
observed to diverge during training; the current PPO-with-clipping
implementation's training stability has not yet been independently
re-validated, and this paper makes no stability claim either way for
the code described in Sections~\ref{subsec3-4}--\ref{subsec3-10}.

Addressing these limitations requires, in order: reconciling the
reward design of Section~\ref{sec4} with the implemented training
reward (or implementing the gap); adding seed control and executing
the Fixed-Time/Max-Pressure/Algorithm evaluation campaign of
Section~\ref{sec5} with multiple seeded runs and variability estimates;
auditing the five networks' signalized-intersection degree distribution
against the four-neighbor attention cap; and, only once an EMV
demand profile exists for at least one scenario, implementing and
evaluating the multi-EMV mechanism of Section~\ref{subsec3-2}.
Extending validation beyond SUMO toward hardware-in-the-loop or
closed-course testing remains a longer-term goal, since SUMO's vehicle
models and signal infrastructure are themselves a simplification of an
operating real-world traffic system. Section~\ref{sec8} discusses these
and related directions in greater detail.


% ===============================================================
% FUTURE WORKS
% ===============================================================

\section{Future Works}\label{sec8}

\subsection{Multi-EMV Routing: From Design to Implementation}

Section~\ref{subsec3-2} sketches, but does not implement, a design in
which each active emergency vehicle contributes a time-decaying bias
term to regional and intersection-level state representations along
its anticipated route, with the Meta-Policy coordinating cross-region
handoff and an arbitration rule resolving conflicts between
simultaneously active EMVs. Realizing this design requires, in order:
(1) adding at least one emergency-vehicle type with origin,
destination, and priority class to a scenario's demand data, since none
of the five scenarios evaluated in this paper currently define one; (2)
implementing the bias-term injection into the Meta-Policy/Sub-Policy
state pipeline; and (3) implementing and evaluating the conflict
arbitration rule under at least two simultaneously active EMVs with
intersecting corridors. Only once these three steps exist does it
become meaningful to compare against EMVLight's dynamic Dijkstra
routing~\cite{emvlight} or to populate Table~\ref{tab:emv-routing}. The
Louvain/Leiden community structure measured in Section~\ref{sec6} may
be a useful abstraction for corridor-level planning once this mechanism
exists, since inter-community transitions correspond to the sparser
boundary links identified by the partitioning step, but this remains an
untested hypothesis rather than a demonstrated capability.

\subsection{Formal Envy-Free Fairness Optimization}

The fairness penalty implemented and evaluated in this paper
(Section~\ref{sec4}: waiting-time variance, a starvation threshold, and
Jain's Fairness Index) does not constitute, and should not be
described as, a formal envy-free guarantee: no envy function comparing
intersections' or vehicles' allocations against one another is defined
or optimized anywhere in the current codebase. A genuine envy-free
extension would need to specify, at minimum, which agents' allocations
are being compared (e.g.\ pairs of neighboring intersections, or
vehicles sharing an approach), a concrete envy function over a
well-defined service quantity, and how a corresponding penalty or
constraint would be incorporated into the PPO objective of
Section~\ref{subsec3-8} without destabilizing training. Only after such
a mechanism is implemented would it become appropriate to re-introduce
``envy-free'' into this framework's name or claims; this paper
deliberately avoids that claim given the current, Jain's-index-based
fairness mechanism alone.

\subsection{Incident and Accident Robustness}

Section~\ref{subsec3-3} proposes, but does not implement, exposing the
controller to injected incidents of varying location, duration, and
severity during training, so that the Sub-Policy learns to detect and
respond to a sudden capacity loss rather than requiring a
hand-engineered override layer. No severity parameter, incident
sampling distribution, or incident-aware training loop currently exists
in the codebase; the only related functionality implemented is a
passive collision counter used for monitoring, not for injecting or
training against incidents. Future work should implement an explicit
incident-injection mechanism on at least one of the five scenarios of
Table~\ref{tab:scenarios}, then evaluate whether incident-aware
training degrades the fairness metrics of Section~\ref{sec4} during
incident-free intervals, as originally intended in
Section~\ref{subsec3-3}. Combining incident robustness, multi-EMV
routing, and the fairness mechanism already implemented represents the
most demanding test of the complete framework and is a plausible target
for the next phase of this research, once each component individually
reaches the implementation status the others already have.


% ===============================================================
% REFERENCES
% ===============================================================

% NOTE (author review required): the entries tagged "% UNVERIFIED" below
% (advancedxlight, tdsp, dynamicrouting, miprouting, rlrouting,
% hierarchicalemv, fairmarl, dtpreemption) carry no author list or
% verifiable venue and could not be confirmed against an external source
% during this revision. They are retained because removing them would
% also require removing their \cite{} occurrences in the text above,
% which this revision did not do without author confirmation. Verify or
% replace before submission.

\begin{thebibliography}{99}

% ---------------------------------------------------------------
% Conventional Traffic Signal Control
% ---------------------------------------------------------------

\bibitem{webster1958}
F.~V. Webster,
``Traffic Signal Settings,''
Road Research Technical Paper No.~39,
Road Research Laboratory, London, UK, 1958.

\bibitem{scoot}
P.~B. Hunt, D.~I. Robertson, R.~D. Bretherton, and R.~D. Winton,
``SCOOT---A Traffic Responsive Method of Coordinating Signals,''
TRRL Laboratory Report 1014,
Transport and Road Research Laboratory, 1981.

\bibitem{scats}
A.~G. Sims and K.~W. Dobinson,
``The Sydney Coordinated Adaptive Traffic System---
Principles, Methodology, Algorithms,''
in \textit{Proceedings of the International Conference on Road
Traffic Signalling}, pp.~67--70, 1980.

\bibitem{opac}
N.~H.~Gartner,
``OPAC: A Demand-Responsive Strategy for Traffic Signal Control,''
\textit{Transportation Research Record},
vol.~906, pp.~75--81, 1983.

\bibitem{maxpressure}
P.~Varaiya,
``Max Pressure Control of a Network of Signalized Intersections,''
\textit{Transportation Research Part C: Emerging Technologies},
vol.~36, pp.~177--195, 2013.

% ---------------------------------------------------------------
% Optimization and Machine Learning
% ---------------------------------------------------------------

\bibitem{genetic}
D.~E. Goldberg,
\textit{Genetic Algorithms in Search, Optimization, and Machine
Learning},
Addison-Wesley, 1989.

\bibitem{fuzzy}
L.~A. Zadeh,
``Fuzzy Logic,''
\textit{Computer},
vol.~21, no.~4, pp.~83--93, 1988.

\bibitem{anntraffic}
M.~J. Dougherty and R.~S. Cobbett,
``A Neural Network Approach to the Prediction of Traffic Flow,''
\textit{Transportation Research Part C},
vol.~2, no.~1, pp.~1--17, 1994.

\bibitem{svmtraffic}
V.~N. Vapnik,
\textit{The Nature of Statistical Learning Theory},
Springer, 1995.

\bibitem{lstmtraffic}
S.~Hochreiter and J.~Schmidhuber,
``Long Short-Term Memory,''
\textit{Neural Computation},
vol.~9, no.~8, pp.~1735--1780, 1997.

% ---------------------------------------------------------------
% Deep Reinforcement Learning for Traffic Signal Control
% ---------------------------------------------------------------

\bibitem{intellight}
H.~Wei, G.~Zheng, H.~Yang, and Z.~Li,
``IntelliLight: A Reinforcement Learning Approach for Intelligent
Traffic Light Control,''
in \textit{Proceedings of the 24th ACM SIGKDD International
Conference on Knowledge Discovery \& Data Mining},
pp.~2496--2505, 2018.

\bibitem{presslight}
H.~Wei, C.~Chen, G.~Zheng, K.~Wu, V.~Gayah, K.~Xu, and Z.~Li,
``PressLight: Learning Max Pressure Control to Coordinate Traffic
Signals in Arterial Network,''
in \textit{Proceedings of the 25th ACM SIGKDD International
Conference on Knowledge Discovery \& Data Mining},
pp.~1290--1298, 2019.

\bibitem{frap}
C.~Zheng, Y.~Xie, H.~Wei, and Z.~Li,
``Learning Phase Competition for Traffic Signal Control,''
in \textit{Proceedings of the 28th International Joint Conference on
Artificial Intelligence}, 2019.

\bibitem{attendlight}
A.~Oroojlooy, M.~Nazari, D.~Hajinezhad, and J.~Silva,
``AttendLight: Universal Attention-Based Reinforcement Learning
Model for Traffic Signal Control,''
in \textit{Advances in Neural Information Processing Systems},
vol.~33, 2020.

\bibitem{metalight}
H.~Wei, G.~Zheng, G.~Yao, and Z.~Li,
``MetaLight: Value-Based Meta-Reinforcement Learning for Traffic
Signal Control,''
in \textit{Proceedings of the AAAI Conference on Artificial
Intelligence}, 2019.

% UNVERIFIED
\bibitem{advancedxlight}
``Advanced-XLight: Advanced Traffic Signal Control with Deep
Reinforcement Learning,''
2020.

% ---------------------------------------------------------------
% Multi-Agent Reinforcement Learning and Graph Methods
% ---------------------------------------------------------------

\bibitem{colight}
H.~Wei, N.~Xu, H.~Zheng, and Z.~Li,
``CoLight: Learning Network-Level Cooperation for Traffic Signal
Control,''
in \textit{Proceedings of the 28th International Joint Conference on
Artificial Intelligence},
pp.~1405--1411, 2019.

\bibitem{mplight}
H.~Wei, G.~Zheng, H.~Yao, and Z.~Li,
``MPLight: Efficient Multi-Agent Reinforcement Learning for
Large-Scale Traffic Signal Control,''
in \textit{Proceedings of the AAAI Conference on Artificial
Intelligence}, 2020.

\bibitem{ma2c}
T.~Chu, J.~Wang, L.~Codec\`a, and Z.~Li,
``Multi-Agent Deep Reinforcement Learning for Large-scale Traffic
Signal Control,''
\textit{IEEE Transactions on Intelligent Transportation Systems},
2019,
doi:10.1109/TITS.2019.2901791.

\url{https://doi.org/10.1109/TITS.2019.2901791}

\bibitem{thousandlights}
C.~Chen, H.~Wei, N.~Xu, G.~Zheng, M.~Yang, Y.~Xiong, K.~Xu,
and Z.~Li,
``Toward A Thousand Lights: Decentralized Deep Reinforcement
Learning for Large-Scale Traffic Signal Control,''
\textit{Proceedings of the AAAI Conference on Artificial
Intelligence},
vol.~34, no.~04, pp.~3414--3421, 2020,
doi:10.1609/aaai.v34i04.5744.

\url{https://doi.org/10.1609/aaai.v34i04.5744}

\bibitem{cityflow}
H.~Zhang, S.~Feng, C.~Liu, Y.~Ding, Y.~Zhu, Z.~Zhou, W.~Zhang,
Y.~Yu, H.~Jin, and Z.~Li,
``CityFlow: A Multi-Agent Reinforcement Learning Environment for
Large Scale City Traffic Scenario,''
in \textit{Proceedings of the World Wide Web Conference},
pp.~3620--3624, 2019,
doi:10.1145/3308558.3314139.

\url{https://doi.org/10.1145/3308558.3314139}

\bibitem{feudalhrl}
F.~Vezhnevets, S.~M.~Osindero, M.~P.~Schaul, R.~H.~Pascanu,
R.~Heess, M.~Jaderberg, M.~W.~Czarnecki, and M.~S.~Rusu,
``FeUdal Networks for Hierarchical Reinforcement Learning,''
in \textit{Proceedings of the 34th International Conference on
Machine Learning},
pp.~3540--3549, 2017.

\bibitem{hirol}
S.~Nachum, S.~Gu, H.~Lee, and S.~Levine,
``Data-Efficient Hierarchical Reinforcement Learning,''
in \textit{Advances in Neural Information Processing Systems},
vol.~31, 2018.

\bibitem{hilight}
Y.~Zhu, H.~Wen, G.~Min, and M.~Luo,
``HiLight: A Hierarchical Reinforcement Learning Framework with
Global Adversarial Guidance for Large-Scale Traffic Signal
Control,''
\textit{arXiv preprint arXiv:2506.14391}, 2025.

\url{https://arxiv.org/abs/2506.14391}

% ---------------------------------------------------------------
% Emergency Vehicle Routing
% ---------------------------------------------------------------

\bibitem{dijkstraemv}
E.~W. Dijkstra,
``A Note on Two Problems in Connexion with Graphs,''
\textit{Numerische Mathematik},
vol.~1, pp.~269--271, 1959.

\bibitem{astar}
P.~E. Hart, N.~J. Nilsson, and B.~Raphael,
``A Formal Basis for the Heuristic Determination of Minimum Cost
Paths,''
\textit{IEEE Transactions on Systems Science and Cybernetics},
vol.~4, no.~2, pp.~100--107, 1968.

% UNVERIFIED
\bibitem{tdsp}
``Time-Dependent Shortest Path Methods for Dynamic Traffic
Routing,''
2020.

% UNVERIFIED
\bibitem{dynamicrouting}
``Dynamic Routing for Emergency Vehicles in Urban Transportation
Networks,''
2020.

% UNVERIFIED
\bibitem{miprouting}
``Mixed Integer Programming Approaches for Emergency Vehicle
Routing,''
2020.

% UNVERIFIED
\bibitem{rlrouting}
``Reinforcement Learning-Based Emergency Vehicle Routing,''
2020.

% ---------------------------------------------------------------
% Emergency Vehicle Routing + Traffic Signal Control
% ---------------------------------------------------------------

\bibitem{emvlight}
H.~Su, Y.~D.~Zhong, J.~Y.~J. Chow, B.~Dey, and L.~Jin,
``EMVLight: A Multi-Agent Reinforcement Learning Framework for an
Emergency Vehicle Decentralized Routing and Traffic Signal Control
System,''
\textit{Transportation Research Part C: Emerging Technologies},
vol.~146, p.~103955, 2023.

\url{https://doi.org/10.1016/j.trc.2022.103955}

% UNVERIFIED
\bibitem{dtpreemption}
H.~Su, H.~Deng, and Y.~Sun,
``Emergency Preemption Without Online Exploration: A Decision
Transformer Approach,''
\textit{arXiv preprint arXiv:2603.22315}, 2026.

\url{https://arxiv.org/abs/2603.22315}

% UNVERIFIED
\bibitem{hierarchicalemv}
``Hierarchical Multi-Agent Actor--Critic Emergency Routing and
Traffic Signal Control,''
2024.

% ---------------------------------------------------------------
% Fairness
% ---------------------------------------------------------------

% UNVERIFIED
\bibitem{fairmarl}
``FairMARL: Fairness-Aware Multi-Agent Reinforcement Learning for
Traffic Signal Control,''
2024.

% ---------------------------------------------------------------
% EMV Visual Sensing
% ---------------------------------------------------------------

\bibitem{nellore2016}
K.~Nellore and G.~P.~Hancke,
``Traffic Management for Emergency Vehicle Priority Based on
Visual Sensing,''
\textit{Sensors},
vol.~16, no.~11, p.~1892, 2016.

\url{https://doi.org/10.3390/s16111892}

% ---------------------------------------------------------------
% Community Detection
% ---------------------------------------------------------------

\bibitem{blondel2008}
V.~D.~Blondel, J.-L.~Guillaume, R.~Lambiotte, and R.~Lefebvre,
``Fast Unfolding of Communities in Large Networks,''
\textit{Journal of Statistical Mechanics: Theory and Experiment},
vol.~2008, no.~10, p.~P10008, 2008.

\bibitem{traag2019}
V.~A.~Traag, L.~Waltman, and N.~J.~van Eck,
``From Louvain to Leiden: Guaranteeing Well-Connected
Communities,''
\textit{Scientific Reports},
vol.~9, no.~1, p.~5233, 2019.

\end{thebibliography}

\end{document}